% =====================================================================
% Слайд 8 — Baseline 3: HSMM (наработка лаборатории)
% =====================================================================
\begin{frame}{Baseline 3: HSMM \small(Cont 2010, эвристика)}
  \footnotesize\setlength{\abovedisplayskip}{2pt}
  \setlength{\belowdisplayskip}{2pt}

  Hidden Semi-Markov Model --- расширение HMM с \rlhi{явной моделью
  длительности}.

  \begin{columns}[T,onlytextwidth]
    \column{0.55\textwidth}
    Канонический вид (Cont 2010):
    \[
      \alpha_t(i) = \sum_{\tau=1}^{D}\sum_{j} \alpha_{t-\tau}(j)\,
        a_{j,i}\, p_i(\tau)\!\!\prod_{k=t-\tau+1}^{t}\!\! b_i(o_k)
    \]
    где $p_i(\tau)$ --- распределение длительности state $i$.

    \vspace{0.4em}
    В нашем коде (\texttt{hsmm\_follower.py}) реализована
    \emph{эвристика:} переходы пересчитываются от
    $r = \texttt{elapsed}/\texttt{nominal\_duration}$. Четыре исхода:
    stay / advance / skip / leap.

    \column{0.45\textwidth}
    \centering
    \resizebox{0.95\linewidth}{!}{%
      % СХЕМА 5 — HSMM как HMM + длительности
\begin{tikzpicture}[
  every node/.style={font=\footnotesize},
  state/.style={draw, thick, circle, minimum size=0.75cm, inner sep=0pt, fill=white},
  arr/.style={-{Latex[length=2mm]}, thick},
]
  % States
  \foreach \i in {0,1,2,3,4} {
    \node[state] (s\i) at (\i*1.4, 0) {$s_{\i}$};
  }

  % advance arrows
  \foreach \i [evaluate=\i as \j using int(\i+1)] in {0,1,2,3} {
    \draw[arr] (s\i) -- (s\j);
  }

  % Tiny p_i(tau) duration curves above each state
  \foreach \i in {0,1,2,3,4} {
    \begin{scope}[shift={(\i*1.4, 1.0)}, scale=0.35]
      \draw[->, thin] (-0.5, 0) -- (0.7, 0)
            node[right, font=\tiny] {$\tau$};
      \draw[->, thin] (-0.5, 0) -- (-0.5, 0.7);
      \draw[thick, smooth, domain=-0.5:0.7, samples=30]
        plot (\x, {0.55*exp(-12*(\x-0.1)^2)});
    \end{scope}
    \node[font=\tiny, above=24pt of s\i] {$p_{\i}(\tau)$};
  }

  % Caption
  \node[align=center, font=\scriptsize\itshape] at (2.8, -1.0)
       {HMM $+$ распределение длительности $p_i(\tau)$\\
        $\Rightarrow$ переходы зависят от \texttt{elapsed} в state};
\end{tikzpicture}
%
    }
  \end{columns}

  \vspace{0.4em}
  {\footnotesize
   \textbf{Добавляет к HMM:} учёт длительностей $\Rightarrow$ правильное
   поведение при rubato и лишних нотах. \quad
   \textbf{Ограничение:} наш HSMM --- эвристика, не канонический Cont.}
\end{frame}
